\chapter{Боль: сигнал тревоги}
\chaptermark{Боль}
\label{sec:pain}
\begin{chaptergoals}
\item почему боль --- приоритетный сигнал тревоги, а не канал измерения;
\item как быстрый и медленный провода сообщают, где повреждение и насколько оно серьёзно;
\item как ворота в спинном мозге дают касанию запретить боль, а широковещательные каналы задают громкость;
\item как сенсибилизация повышает усиление после повреждения и как боль учит и прерывает работу.
\end{chaptergoals}

\section{Тревога, а не измерение}

Все датчики главы~\ref{sec:sensors} измеряли мир: сколько света, какой высоты звук, какое давление. Боль устроена иначе. Она сообщает не «что там», а «опасно, брось всё». Для инженера это не измерительный канал, а \emph{сигнал тревоги}: линия с высшим приоритетом, которая должна пробиться сквозь любую текущую работу машины.

От измерительного канала сигнализация отличается тремя свойствами, и все три у боли есть. Во-первых, её трудно приглушить привыканием: датчики боли привыкают гораздо слабее датчиков касания, а после лёгкого повреждения становятся даже чувствительнее~\cite{LaMotte1982hyper}; ослабление ответа после сильного нагрева измерено лишь у одного их типа~\cite{Treede1995heat}. Датчик света на постоянном фоне замолкает (рис.~\ref{fig:adaptation}), а сигнализация, которая замолчала через минуту, бесполезна. Во-вторых, у неё высокий порог: датчики боли срабатывают только тогда, когда воздействие близко к опасному, например на нагрев --- у разных датчиков пороги от $41^\circ$ до $46$--$48^\circ$, в зависимости от типа~\cite{Treede1995heat, Weidner1999cfib}. В-третьих --- и это главное для этой главы --- громкость тревоги решает не только датчик, но и сама машина. Одна и та же рана болит по-разному в разной обстановке. Посмотрим, из каких уже знакомых деталей это собрано.

Датчики боли --- \nt{GLUT}-нейроны, как и остальные чувствительные нейроны главы~\ref{sec:sensors}. Часть из них вместе с глутаматом выбрасывает вещество~P из приложения~\ref{sec:othermediators}, медленно усиливающее передачу.

\section{Два провода: быстрая и медленная боль}

Ударьте палец --- и вы почувствуете боль дважды: сначала короткую и резкую, а через секунду тупую и долгую. Это два разных провода. Быстрый провод изолирован и проводит импульсы со скоростью $v$ от $5$ до $30$ метров в секунду; медленный тонкий и голый, $0.5$--$2$ метра в секунду. Сигнал от ступни проходит около метра, и время прихода
\begin{empheq}[box=\keybox]{equation}
t \;=\; \frac{L}{v}
\label{eq:paintime}
\end{empheq}
для быстрого провода --- десятки миллисекунд, для медленного --- около секунды. Два провода несут два сообщения. Быстрый говорит \emph{где} и \emph{когда}: его импульсы приходят раньше и точнее, как в коде временем главы~\ref{sec:code}. Медленный говорит \emph{насколько плохо}, и его долгая тупая боль держит машину в режиме тревоги, пока повреждение не прошло.

\section{Ворота в спинном мозге}

Первое, куда приходит сигнал боли, --- спинной мозг, и там он проходит через \emph{ворота}~\cite{MelzackWall1965}; конкретные нейроны этих ворот найдены сравнительно недавно~\cite{Duan2014}. На выходной \nt{GLUT}-нейрон, который передаёт боль дальше в мозг, сходятся провод боли и провод прикосновения. А рядом сидит тормозящий нейрон-посредник, \nt{GABA} или \nt{GLYC}, который прикосновение возбуждает (рис.~\ref{fig:gate}). Прикосновение закрывает ворота. В исходной теории ворот~\cite{MelzackWall1965} боль, наоборот, выключает этого посредника, и сильная боль ворота открывает; это допущение теории, на найденных нейронах ворот оно прямо не показано.

\begin{figure}[t]
\centering
\begin{tikzpicture}[>=Stealth,
  nrn/.style={circle,draw,very thick,minimum size=1.0cm,inner sep=0pt,font=\scriptsize},
  inh/.style={-{Bar[width=7pt]},thick,red!70!black},
  bc/.style={->,thick,densely dashed,orange!85!black}]
\node[nrn,fill=blue!6] (Z) at (6,0) {\nt{GLUT}};
\node[nrn,fill=red!10] (Q) at (3.6,-1.6) {\nt{GABA}};
\draw[->,very thick,red!70!black] (0,0.6) node[left,font=\small,black,align=right] {боль $\nu$} -- (5.45,0.15);
\draw[->,very thick,blue!60!black] (0,-2.6) node[left,font=\small,black,align=right] {касание $\mu$} -- (2.9,-2.0);
\draw[->,thick,blue!60!black] (1.8,-2.35) -- (5.5,-0.35) node[pos=0.8,below right=-2pt,font=\scriptsize] {$w_{z\mu}$};
\draw[inh] (1.6,0.5) -- (3.35,-1.1) node[pos=0.5,left=1pt,font=\scriptsize] {$w_{q\nu}$};
\draw[inh] (Q) -- (Z) node[pos=0.5,above left=-1pt,font=\scriptsize] {$\beta$};
\draw[->,very thick] (Z) -- (8.2,0) node[right,font=\small,align=left] {в мозг: $z$};
\draw[bc] (6,2.1) node[above,font=\scriptsize,black,align=center] {сверху: \nt{SERO}, \nt{NORA}, опиоиды} -- (Z.north) node[pos=0.45,right,font=\scriptsize,black] {усиление $g$};
\end{tikzpicture}
\caption{Ворота боли в спинном мозге по модели~\eqref{eq:gate}. Выходной \nt{GLUT}-нейрон получает боль $\nu$ и слабое возбуждение от касания $\mu$. Тормозящий посредник (\nt{GABA} или \nt{GLYC}) возбуждается касанием и, по допущению теории ворот, выключается болью. Сверху, по широковещательным каналам, приходит усиление $g$.}
\label{fig:gate}
\end{figure}

Запишем это для частот, как в главе~\ref{sec:gaba}. Частота посредника $q$ и выходная частота $z$ равны
\begin{empheq}[box=\keybox]{equation}
q \;=\; \bigl[w_{q\mu}\,\mu + q_0 - w_{q\nu}\,\nu\bigr]_+ ,
\qquad
z \;=\; \bigl[\,g\,(\nu + w_{z\mu}\,\mu) - \beta\,q\,\bigr]_+ ,
\label{eq:gate}
\end{empheq}
где $\nu$ --- вход боли, $\mu$ --- вход касания, $w$ --- веса связей (первый индекс --- кому, второй --- от кого), $\beta$ --- сила торможения, $q_0$ --- собственная фоновая активность посредника, $g$ --- усиление, которое задают сверху (о нём ниже). Это запрет~\eqref{eq:veto} из главы~\ref{sec:gaba}: «боль, но не при касании», --- с одной поправкой, взятой из теории ворот как допущение: сильная боль сама выключает запрещающий нейрон.

Мы посчитали модель при $w_{q\mu}=1$, $q_0=0.2$, $w_{q\nu}=0.5$, $w_{z\mu}=0.3$, $\beta=1.5$ (рис.~\ref{fig:gatecurves}). Без касания боль начинает проходить с $\nu\approx0.18$. С касанием порог поднимается до $0.86$, и при $\nu=1$ выход падает вчетверо, с $1$ до $0.25$. А при сильной боли, $\nu=2$, касание уже ничего не меняет: ворота распахнуты. Так и в жизни: потереть ушиб помогает, а при серьёзной травме --- нет. Само касание, без боли, через ворота не проходит: тревога не поднимается от прикосновения.

\begin{figure}[t]
\centering
\begin{tikzpicture}[>=Stealth,x=5.2cm,y=1.35cm]
\draw[->,gray!70!black] (0,0) -- (2.12,0) node[right,font=\small,black] {боль $\nu$};
\draw[->,gray!70!black] (0,0) -- (0,4.3) node[left,font=\small,black] {$z$};
\foreach \t in {0,0.5,1,1.5,2}{\draw[gray!60!black] (\t,-0.05) -- (\t,0.05); \node[below,font=\scriptsize] at (\t,0) {\t};}
\foreach \f in {1,2,3,4}{\draw[gray!60!black] (-0.01,\f) -- (0.01,\f); \node[left,font=\scriptsize] at (0,\f) {\f};}
\draw[very thick,red!70!black] plot file {figures/pain_gate_notouch.dat};
\draw[very thick,blue!60!black] plot file {figures/pain_gate_touch.dat};
\draw[very thick,orange!85!black,densely dashed] plot file {figures/pain_gate_sens.dat};
\draw[very thick,red!70!black] (0.08,3.9) -- (0.2,3.9) node[right,font=\scriptsize,black] {без касания};
\draw[very thick,blue!60!black] (0.08,3.45) -- (0.2,3.45) node[right,font=\scriptsize,black] {с касанием};
\draw[very thick,orange!85!black,densely dashed] (0.08,3.0) -- (0.2,3.0) node[right,font=\scriptsize,black] {после сенсибилизации, $g=2$};
\end{tikzpicture}
\caption{Выход ворот~\eqref{eq:gate} в зависимости от силы боли. Касание поднимает порог и приглушает слабую и среднюю боль, но не сильную. Сенсибилизация (штриховая линия) вдвое увеличивает усиление: та же боль передаётся вдвое громче, а порог опускается.}
\label{fig:gatecurves}
\end{figure}

\section{Ручка громкости сверху}

Ворота управляются не только снизу. Из ствола мозга к ним идут нисходящие пути, которые меняют усиление $g$ в~\eqref{eq:gate}: \nt{SERO}, \nt{NORA} и опиоидные пептиды из приложения~\ref{sec:othermediators}~\cite{Fields2004}. Это те же широковещательные каналы, что и в главах~\ref{sec:serocomputer}--\ref{sec:acet}, только здесь их ручка --- громкость тревоги. Они умеют крутить её в обе стороны: одни и те же нисходящие цепи могут и приглушать боль, и усиливать её.

Зачем машине убавлять тревогу? Затем, что тревога --- это прерывание, а прерывание не всегда уместно. Животное, убегающее от хищника, не должно останавливаться из-за раненой лапы: сначала бежать, потом зализывать рану. Ожидание облегчения тоже приглушает боль, и часть этого эффекта снимается, если заблокировать опиоидные рецепторы~\cite{Levine1978}: ожидание, вычисленное в мозге, спускается по опиоидному каналу к воротам. Сама по себе эта картина измерена; как именно мозг решает, какой должна быть громкость, --- гипотеза.

\section{Сенсибилизация: тревога, которая учится}

После повреждения выходные нейроны ворот становятся чувствительнее: тот же вход вызывает больший выход, а порог опускается~\cite{Woolf1983}. В модели~\eqref{eq:gate} это рост усиления $g$, и простейшее его описание --- медленная RC-цепь, которую накачивает сама боль:
\begin{empheq}[box=\keybox]{equation}
\tau_s\,\frac{\dd g}{\dd t} \;=\; -(g-1) + k_s\,z ,
\label{eq:sensitization}
\end{empheq}
где $\tau_s$ --- время, за которое чувствительность возвращается к норме; оно больше длительности самой боли, потому что сенсибилизация держится и после того, как повреждающий стимул прекратился~\cite{Woolf1983}, а $k_s$ --- сила накачки. Пока ткань повреждена, выход ворот держит $g$ выше единицы, и всё, что происходит рядом с раной, передаётся громче (штриховая линия на рис.~\ref{fig:gatecurves}: при $g=2$ порог опускается до $0.11$, а выход удваивается). Это разумная настройка сигнализации: пока повреждение не зажило, лучше перестраховаться.

Для инженера это положительная обратная связь с медленной утечкой: боль повышает усиление, усиление --- боль. Пока петля слабая, $g$ возвращается к норме, когда рана зажила. Если же петля сильная, сигнализация может застрять во включённом состоянии и после того, как повреждения уже нет, --- как группа с сильной обратной связью на рис.~\ref{fig:bistable}. Модель~\eqref{eq:sensitization} --- упрощение; то, что центральная сенсибилизация существует, измерено.

\section{Боль как учитель и как прерывание}

У боли в машине две работы помимо сигнализации.

\emph{Учитель.} Боль --- отрицательная награда: в формуле ошибки~\eqref{eq:ctd} она входит как $r<0$. Всё, что говорилось в главе~\ref{sec:dofa}, работает и здесь: предвестник боли начинает вызывать ошибку заранее, и сеть учится избегать того, что к боли ведёт. Опыты на людях показывают, что активность мозга при обучении на боли следует именно ошибке временных разностей~\cite{Seymour2004}, а часть \nt{DOPA}-нейронов отвечает и на неприятные события (глава~\ref{sec:dofaalt}).

\emph{Прерывание.} Боль отрывает внимание от текущей задачи --- это её назначение, а не побочный эффект~\cite{EcclestonCrombez1999}. В главе~\ref{sec:nora} та же роль «прерывания» обсуждалась для отрывистой пачки \nt{NORA}. А самый быстрый ответ даже не доходит до мозга: отдёргивание руки от горячего замыкается прямо в спинном мозге той же парой мышц и тем же запретом противника, что на рис.~\ref{fig:antagonists}.

\begin{keyblock}
\begin{proposition}[Сигнал тревоги]
\label{prop:pain}
Боль --- не измерение, а сигнал тревоги с высоким приоритетом. Она привыкает гораздо слабее измерительных каналов, идёт по двум проводам (быстрому --- «где», медленному --- «насколько плохо»), проходит через ворота, которые закрывает касание (а сильная боль, по допущению теории ворот, открывает), и её громкость задают сверху широковещательные каналы. Эти устройства измерены, кроме открывания ворот болью; простые модели~\eqref{eq:gate} и~\eqref{eq:sensitization} --- упрощения.
\end{proposition}
\end{keyblock}

\section{Итог}

\begin{table}[t]
\centering
\footnotesize
\setlength{\tabcolsep}{4pt}
\renewcommand{\arraystretch}{1.15}
\begin{tabular}{>{\raggedright}p{3.0cm}>{\raggedright}p{4.9cm}>{\raggedright\arraybackslash}p{4.9cm}}
\toprule
Что делает боль & Как & Что известно \\
\midrule
поднимает тревогу & высокий порог, привыкание слабее, чем у касания & порог измерен; сравнение привыкания --- оценка \\
сообщает «где» и «насколько плохо» & два провода с разной скоростью~\eqref{eq:paintime} & измерено \\
проходит через ворота & запрет через \nt{GABA}/\nt{GLYC}~\eqref{eq:gate} & ворота измерены; открывание болью --- допущение теории; модель --- упрощение \\
меняет громкость & нисходящие \nt{SERO}, \nt{NORA}, опиоиды & измерено; правило выбора громкости --- гипотеза \\
учится после повреждения & рост усиления~\eqref{eq:sensitization} & сенсибилизация измерена; модель --- упрощение \\
учит избегать & отрицательная награда в~\eqref{eq:ctd} & сходство с ошибкой временных разностей измерено \\
прерывает работу & захват внимания; рефлекс отдёргивания & измерено \\
\bottomrule
\end{tabular}
\caption{Боль как сигнал тревоги.}
\label{tab:pain}
\end{table}

Боль собрана из деталей, которые мы уже видели (таблица~\ref{tab:pain}): \nt{GLUT}-датчики, два провода, запрет через тормозящего посредника, широковещательная ручка громкости, медленная RC-цепь, ошибка предсказания и прерывание. Новое в ней одно: это единственный вход машины, громкость которого машина задаёт сама, --- сигнализация, которую владелец может и приглушить, и перенастроить.

\section{Задачи}

\begin{exercise}[\lvlA]
\label{ex:14:1}
\emph{Два провода.} Сигнал идёт от ступни, $L=1$~м. (а)~Залп идёт по пучку проводов одного типа, скорости которых заполняют диапазон~\eqref{eq:paintime}. На сколько он растягивается к приходу в спинной мозг для быстрого и медленного типа? (б)~Волны боли по проводам $15$ и $1$~м/с различимы при разнице от $0.1$~с. С какого расстояния они сливаются?
\end{exercise}

\begin{exercise}[\lvlB]
\label{ex:14:2}
\emph{Когда касание болит.} Ворота~\eqref{eq:gate} с параметрами текста, усиление $g$ растёт при сенсибилизации. До какого $g$ касание без боли не проходит через ворота ни при какой силе $\mu$? При каком $g$ начинает проходить касание $\mu=1$?
\end{exercise}

\begin{exercise}[\lvlB]
\label{ex:14:3}
\emph{Устойчивость сенсибилизации.} Входы постоянны, выход ворот линеен по $g$: $z=gA-C$, $A=\nu+w_{z\mu}\mu$, $C=\beta q\ge0$. (а)~Найдите равновесие $g^*$ модели~\eqref{eq:sensitization} и условие его устойчивости. (б)~При $k_s=0.5$ найдите силу боли, выше которой модель уходит вразнос, без касания и с касанием $\mu=1$.
\end{exercise}

\begin{exercise}[\lvlA]
\label{ex:14:4}
\emph{Предвестник боли.} Сигнал в момент $0$ предсказывает боль в момент $T=2$~с, в~\eqref{eq:ctd} $r(t)=-P\,\delta(t-T)$, $\tau_\gamma=2$~с. (а)~Найдите площадь всплеска $\delta$ в момент сигнала. (б)~Действие в момент $t_e=1$~с отменяет боль. Каков всплеск $\delta$ в этот момент и чему он научит сеть?
\end{exercise}

\begin{exercise}[\lvlB]
\label{ex:14:5}
\emph{Сигнализация, которая защёлкивается.} Дополните ворота из \texttt{models/\allowbreak pain\_gate.py} динамикой~\eqref{eq:sensitization} и насыщением $z\le4$. Касание $\mu=1$ есть всегда, повреждение $\nu=2$ длится время $T$, затем $\nu=0$. Моделированием найдите наименьшее $T$ (в единицах $\tau_s$), после которого сигнализация остаётся включённой, при $k_s=4$, $4.6$, $5$, $6$, $8$.
\end{exercise}

\begin{exercise}[\lvlB]
\label{ex:14:6}
\emph{Проектирование ворот.} При $\beta=1.5$, $w_{z\mu}=0.3$, $g=1$ подберите в~\eqref{eq:gate} $q_0$, $w_{q\nu}$, $w_{q\mu}$ так, чтобы порог боли был $0.2$ без касания и $1.0$ с касанием $\mu=1$, а касание переставало приглушать боль при $\nu_\times=2.5$. До какого усиления $g$ касание без боли не проходит через ворота?
\end{exercise}

\begin{exercise}[\lvlC]
\label{ex:14:7}
\emph{Ручка громкости.} Работа приносит ценность $V$ в единицу времени, работа с повреждением $\nu$ стоит $c\nu$, так что прерывать выгодно при $\nu>V/c$; боль прерывает при $z>\theta=0.5$. (а)~Какое усиление $g^*$ ворот~\eqref{eq:gate} без касания даёт этот порог при $V/c=0.25$, $0.5$, $2$? (б)~Из каких сигналов книги машина могла бы оценивать $V$ и $c$?
\end{exercise}
